% ============================================================================
% COURS 13 — EXERCICES, PARTIE A
% Source : 17-MAS_MS.docx
% ============================================================================

\section{Exercices du chapitre}
\label{sec:masms-8}

\subsection{Banc BALAFRE}

\textit{Source : concours ATS 2019.}

\begin{center}
\TLMAInlineImage[.28\linewidth]{image79.png}\hfill
\TLMAInlineImage[.62\linewidth]{image80.png}
\end{center}

Le banc BALAFRE (BAnc d'essais à LAmes Fluides à haut REynolds) permet d'identifier le comportement dynamique de joints annulaires utilisés dans les turbomachines. Un moteur asynchrone triphasé entraîne le rotor du banc. Pendant l'essai, sa vitesse doit être stabilisée à $6000\ \mathrm{tr\,min^{-1}}$ malgré un couple résistant estimé à $C_{res}=300\ \mathrm{N\,m}$.

Le réseau fournit $230/400\ \mathrm V$ à $50\ \mathrm{Hz}$. La machine est commandée par un variateur fonctionnant à $U_s/f$ constant. Les pertes Joule statoriques, les pertes fer et les pertes mécaniques sont négligées.

\TLMAImage[.45\linewidth]{image81.png}

\begin{enumerate}
  \item À partir de la plaque signalétique, déterminer le couplage, le nombre de paires de pôles, la vitesse de synchronisme, le glissement nominal et le couple utile nominal.
  \item Le modèle équivalent d'une phase comporte une inductance de magnétisation $L_0$, une inductance de fuite totale $L_c$ et une résistance rotorique ramenée au stator $R/g$. Établir :
  \[
  C_u=\frac{3pU_s^2}{\omega}\,
  \frac{R/g}{(R/g)^2+(L_c\omega)^2}.
  \]
  \item Identifier sur la caractéristique couple--vitesse le démarrage, le point nominal, le synchronisme et la zone instable.
  \item Déterminer le glissement $g_M$ correspondant au couple maximal et exprimer $C_M$.
  \item Le constructeur indique $C_M/C_N=3{,}5$. Déterminer $L_c$, puis estimer $R$ à partir du point nominal.
  \item À l'aide des courbes fournies pour des fréquences comprises entre $90$ et $110\ \mathrm{Hz}$, déterminer la fréquence à imposer pour maintenir $6000\ \mathrm{tr\,min^{-1}}$ sous $300\ \mathrm{N\,m}$.
  \item Vérifier que la durée d'accélération imposée par le cahier des charges reste inférieure à $5\ \mathrm s$.
\end{enumerate}

\subsection{Panneaux publicitaires déroulants}

\textit{Source : concours ATS 2011.}

\begin{minipage}{.73\linewidth}
Le panneau déroulant de type Senior utilise deux moteurs asynchrones commandés par des variateurs scalaires. Les affiches défilent à $1\ \mathrm{m\,s^{-1}}$ avec des rampes d'accélération et de décélération d'une seconde.

En régime permanent, le couple moteur nécessaire vaut $C_m=0{,}28\ \mathrm{N\,m}$ pour une vitesse $\Omega_m=139\ \mathrm{rad\,s^{-1}}$. Le motoréducteur possède les caractéristiques nominales suivantes sous $50\ \mathrm{Hz}$ :
\[
P_N=120\ \mathrm W,\quad N_N=1300\ \mathrm{tr\,min^{-1}},\quad
C_N=0{,}88\ \mathrm{N\,m},\quad C_{max}/C_N=1{,}9.
\]
\end{minipage}\hfill
\begin{minipage}{.23\linewidth}
\TLMAInlineImage[\linewidth]{image86.jpeg}
\end{minipage}

Le couple est modélisé par :
\[
C_m=\frac{2C_{max}}{g_0/g+g/g_0}.
\]

\begin{enumerate}
  \item Calculer la puissance mécanique nécessaire et montrer que le moteur est surdimensionné.
  \item Déterminer $p$, $\Omega_s$, $N_s$ et le glissement nominal.
  \item Calculer $g_0$, puis tracer $C_m(g)$ pour $0<g<1$ en repérant la zone stable.
  \item Pour le couple réel $C_{maff}=0{,}30\ \mathrm{N\,m}$, utiliser l'approximation de faible glissement pour déterminer $g_{aff}$.
  \item Montrer que, dans cette zone, $C_m=\lambda(\Omega_s-\Omega_m)$ et préciser $\lambda$.
  \item La commande est réalisée à $U/f$ constant et l'on prendra $\lambda=0{,}0453\ \mathrm{N\,m\,s\,rad^{-1}}$. Déterminer la vitesse de synchronisme puis la fréquence à programmer pour obtenir $\Omega_m=139\ \mathrm{rad\,s^{-1}}$.
\end{enumerate}

\subsection{Entraînement de la broche d'un centre d'usinage}

\textit{Source : concours CCP TSI 2002.}

\begin{minipage}{.72\linewidth}
Le centre d'usinage horizontal étudié possède une broche entraînée par une machine asynchrone triphasée et un onduleur de tension. Les caractéristiques nominales sous $50\ \mathrm{Hz}$ sont :
\[
P_N=5{,}35\ \mathrm{kW},\quad C_N=35\ \mathrm{N\,m},\quad
N_N=1460\ \mathrm{tr\,min^{-1}},\quad I_N=9{,}5\ \mathrm A,
\]
avec $C_D/C_N=2{,}4$ et $C_{max}=89{,}3\ \mathrm{N\,m}$.
\end{minipage}\hfill
\begin{minipage}{.24\linewidth}
\TLMAInlineImage[\linewidth]{image88.jpeg}
\end{minipage}

Toutes les pertes sont négligées sauf les pertes Joule rotoriques. Le moteur fonctionne à couple maximal constant jusqu'à $1500\ \mathrm{tr\,min^{-1}}$, puis à puissance utile maximale constante jusqu'à $6000\ \mathrm{tr\,min^{-1}}$.

\begin{enumerate}
  \item Déterminer $p$, $N_s$, $g_N$, $P_{tr}$, les pertes Joule rotoriques, le rendement et le facteur de puissance au point nominal.
  \item Tracer l'enveloppe couple--vitesse autorisée du moteur.
  \item Reporter sur le même graphe les couples résistants ramenés au moteur pour les deux rapports de boîte :
  \[
  C_{r1}=12+0{,}15\Omega,
  \qquad
  C_{r2}=1+0{,}0125\Omega.
  \]
  \item Déterminer les points de fonctionnement et les couples accélérateurs disponibles.
  \item L'onduleur est commandé à $U/f$ constant et délivre $400\ \mathrm V$ entre phases à $50\ \mathrm{Hz}$. Tracer le coefficient de modulation en fonction de la fréquence jusqu'à $80\ \mathrm{Hz}$.
  \item À partir du schéma équivalent simplifié, établir l'expression de $C_{em}$, calculer $C_{max}$ à $50\ \mathrm{Hz}$ et rechercher la fréquence donnant le couple accélérateur maximal au démarrage.
\end{enumerate}

\subsection{Maison hantée}

\textit{Source : concours ATS 2009.}

\begin{center}
\TLMAInlineImage[.28\linewidth]{image92.png}\hfill
\TLMAInlineImage[.62\linewidth]{image93.jpeg}
\end{center}

Seize véhicules circulent sur une piste guidée. Chaque véhicule est entraîné par un moteur asynchrone triphasé associé à un variateur. La vitesse de translation doit rester comprise entre $1$ et $1{,}3\ \mathrm{m\,s^{-1}}$, avec une variation maximale tolérée de $10\,\%$. Le moteur SEW DV100M4, connecté en triangle, possède sous $50\ \mathrm{Hz}$ :
\[
P_N=2{,}2\ \mathrm{kW},\quad C_N=15\ \mathrm{N\,m},\quad
N_N=1410\ \mathrm{tr\,min^{-1}},\quad C_{max}=40\ \mathrm{N\,m}.
\]

Le couple électromagnétique est modélisé par :
\[
C_{em}=\frac{3U^2}{\Omega_s}\frac{R_2g}{R_2^2+(X_2g)^2}.
\]

\begin{enumerate}
  \item Déterminer $g_M$ et $C_M$, puis écrire la relation réduite :
  \[
  C_{em}=2C_M\frac{x}{1+x^2},\qquad x=\frac{g}{g_M}.
  \]
  \item Déterminer $p$ et le glissement nominal. Utiliser les données constructeur pour calculer $g_M$.
  \item Linéariser la caractéristique autour du synchronisme et exprimer $C_{em}$ en fonction de $f$, $p$ et $\Omega_m$.
  \item Déterminer la fréquence qui assure $1{,}3\ \mathrm{m\,s^{-1}}$ pour un couple moteur de $10\ \mathrm{N\,m}$.
  \item Étudier la réversibilité du convertisseur et préciser le trajet de l'énergie lorsque le couple devient négatif.
  \item Calculer la vitesse lorsque le couple passe de $10$ à $-4\ \mathrm{N\,m}$ et conclure sur la nécessité d'une régulation de vitesse.
\end{enumerate}
